% MANUAL DE USUARIO - Super Asgard-V
%
% Documento principal: portada + control de versiones + índices.
% Compilar con pdflatex DOS veces: la portada usa "remember picture" y el índice se genera a partir del .toc de la pasada anterior.
\documentclass[12pt,a4paper,oneside,openany]{book}

\usepackage[utf8]{inputenc}
\usepackage[spanish,es-tabla]{babel}
\usepackage[T1]{fontenc}
\usepackage{lmodern}
\usepackage[papersize={210mm,297mm},tmargin=25mm,bmargin=25mm,lmargin=20mm,rmargin=20mm]{geometry}
\setlength{\headheight}{16pt}

% No indentar los párrafos y separarlos 1em
\setlength\parindent{0cm}
\setlength\parskip{1em}
\usepackage{tikz}
\usetikzlibrary{calc}

% --- Paquetes que necesita manual.tex
\usepackage{subfiles}
\usepackage{graphicx}
\usepackage{float}
\usepackage{listings}
\usepackage{xcolor}
\usepackage{fancyhdr}
\usepackage{titlesec}
\usepackage{booktabs}
\usepackage{tabularx}
\usepackage{hyperref}

% Títulos de capítulo sin el espacio enorme por defecto
\titleformat{\chapter}[display]
{\normalfont\huge\bfseries}{\chaptertitlename\ \thechapter}{20pt}{\Huge}
\titlespacing*{\chapter}{0pt}{-50pt}{40pt}

% \paragraph como título de bloque: el contenido empieza en la línea siguiente
\setcounter{secnumdepth}{4}
\titleformat{\paragraph}{\normalfont\normalsize\bfseries}{\theparagraph}{1em}{}
\titlespacing*{\paragraph}{0pt}{3.25ex plus 1ex minus .2ex}{1.5ex plus .2ex}

% Los espacios sobrantes van al final de página, no entre el texto
\raggedbottom

% Rutas donde buscar las capturas del manual
\graphicspath{{./}{./img/}{./imagenes/}{./figuras/}}

\hypersetup{
  colorlinks=true,
  linkcolor=black,
  urlcolor=blue,
  pdftitle={Super Asgard-V: Manual de Usuario},
  pdfauthor={Jorge Márquez Torres}
}

% Listados de código
\addto\captionsspanish{
  \renewcommand{\lstlistingname}{Listado}
  \renewcommand{\lstlistlistingname}{Índice de listados}
}
\definecolor{codegreen}{rgb}{0,0.6,0}
\definecolor{codegray}{rgb}{0.5,0.5,0.5}
\definecolor{codepurple}{rgb}{0.58,0,0.82}
\definecolor{backcolour}{rgb}{0.95,0.95,0.92}
\lstdefinestyle{asgardstyle}{
  commentstyle=\color{codegreen},
  keywordstyle=\color{blue},
  stringstyle=\color{codepurple},
  backgroundcolor=\color{backcolour},
  basicstyle=\footnotesize\ttfamily\color{black},
  belowcaptionskip=1\baselineskip,
  breakatwhitespace=false,
  breaklines=true,
  captionpos=b,
  escapeinside={\%*}{*)},
  firstnumber=1,
  frame=single,
  keepspaces=true,
  numbers=left,
  numbersep=10pt,
  numberstyle=\tiny\color{gray},
  rulecolor=\color{black},
  showspaces=false,
  showstringspaces=false,
  showtabs=false,
  stepnumber=1,
  tabsize=2,
  title=\lstname,
  xleftmargin=1.5cm,
  % Tildes y eñes en los listados (listings no soporta UTF-8 directamente)
  literate={á}{{\'a}}1 {é}{{\'e}}1 {í}{{\'i}}1 {ó}{{\'o}}1 {ú}{{\'u}}1
           {Á}{{\'A}}1 {É}{{\'E}}1 {Í}{{\'I}}1 {Ó}{{\'O}}1 {Ú}{{\'U}}1
           {ñ}{{\~n}}1 {Ñ}{{\~N}}1
}
\lstset{style=asgardstyle}

% Cabeceras: número de página y capítulo actual
\pagestyle{fancy}
\fancyhf{}
\fancyhead[L]{\leftmark}
\fancyfoot[C]{\thepage}

% Páginas de inicio de capítulo: sin cabecera, número abajo centrado
\fancypagestyle{plain}{%
  \fancyhf{}%
  \fancyfoot[C]{\thepage}%
  \renewcommand{\headrulewidth}{0pt}%
}

% --- Datos de la portada
\newcommand{\portadaTitulo}{Super Asgard-V}
\newcommand{\portadaSubtitulo}{Un simulador moderno para aprender arquitectura de computadores, \\directamente en el navegador}
\newcommand{\portadaDocumento}{Manual de Usuario}
\newcommand{\portadaAutor}{Jorge Márquez Torres}
\newcommand{\portadaDirector}{José Sánchez Moreno}
\newcommand{\portadaFecha}{2026}

% --- Colores (los mismos que usa la aplicación)
\definecolor{asgNavy}{HTML}{0F172A}
\definecolor{asgNavyLight}{HTML}{1E3A5F}
\definecolor{asgAccent}{HTML}{3B82F6}
\definecolor{asgSlate}{HTML}{94A3B8}
\definecolor{asgFlame}{HTML}{F97316}
\definecolor{asgFlameCore}{HTML}{FACC15}

% --- Logo: cohete (inspirado en el icono "rocket_launch" de la app)
\newcommand{\asgRocket}{%
  % Llama
  \fill[asgFlame] (-0.45,-1.55) .. controls (-0.55,-2.3) and (-0.1,-2.8) .. (0,-3.3)
                  .. controls (0.1,-2.8) and (0.55,-2.3) .. (0.45,-1.55) -- cycle;
  \fill[asgFlameCore] (-0.25,-1.55) .. controls (-0.3,-2.0) and (-0.05,-2.35) .. (0,-2.6)
                      .. controls (0.05,-2.35) and (0.3,-2.0) .. (0.25,-1.55) -- cycle;
  % Aletas
  \fill[asgAccent] (-0.8,-0.1) -- (-1.65,-0.95) -- (-1.65,-1.85) -- (-0.8,-1.2) -- cycle;
  \fill[asgAccent] ( 0.8,-0.1) -- ( 1.65,-0.95) -- ( 1.65,-1.85) -- ( 0.8,-1.2) -- cycle;
  % Tobera
  \fill[asgSlate] (-0.55,-1.2) rectangle (0.55,-1.55);
  % Cuerpo
  \fill[white] (-0.8,-1.2) -- (-0.8,1.0)
               .. controls (-0.8,2.2) and (-0.3,3.0) .. (0,3.45)
               .. controls (0.3,3.0) and (0.8,2.2) .. (0.8,1.0)
               -- (0.8,-1.2) -- cycle;
  % Aleta central
  \fill[asgAccent] (-0.12,-0.4) rectangle (0.12,-1.2);
  % Ventana
  \fill[asgAccent] (0,1.3) circle (0.42);
  \fill[asgNavy]   (0,1.3) circle (0.28);
  \fill[white, opacity=0.8] (-0.08,1.4) circle (0.07);
}

\begin{document}

\begin{titlepage}
\begin{tikzpicture}[remember picture, overlay]
  % Coordenadas relativas a la página
  \coordinate (NO) at (current page.north west);
  \coordinate (SE) at (current page.south east);
  \coordinate (corte) at ($(current page.north)!0.62!(current page.south)$);

  % Zona superior oscura
  \shade[top color=asgNavy, bottom color=asgNavyLight]
    (NO) rectangle (SE |- corte);

  % Estrellas
  \foreach \x/\y/\r in {
      1.6/-2.1/0.035, 3.4/-4.3/0.025, 5.2/-1.5/0.03, 2.3/-7.8/0.03,
      4.1/-10.9/0.025, 1.3/-12.6/0.035, 6.6/-6.2/0.02, 7.9/-2.8/0.03,
      13.4/-1.9/0.025, 15.2/-4.6/0.035, 17.8/-2.4/0.03, 19.4/-5.8/0.025,
      18.6/-9.4/0.035, 16.3/-11.8/0.025, 19.8/-13.1/0.03, 12.2/-12.9/0.02,
      9.4/-1.1/0.02, 11.1/-3.6/0.025, 14.6/-8.1/0.02, 3.0/-15.4/0.025,
      18.1/-16.2/0.03, 7.2/-14.8/0.02} {
    \fill[white, opacity=0.7] ($(NO)+(\x,\y)$) circle (\r);
  }

  % Halo y logo
  \coordinate (logo) at ($(current page.north)+(0,-6.6)$);
  \fill[asgAccent, opacity=0.10] (logo) circle (3.9);
  \fill[asgAccent, opacity=0.14] (logo) circle (2.9);
  \begin{scope}[shift={(logo)}, rotate=-40, scale=1.05]
    \asgRocket
  \end{scope}

  % Título y subtítulo
  \node[anchor=north, text=white, font=\sffamily\bfseries\fontsize{46}{50}\selectfont]
    at ($(current page.north)+(0,-11.4)$) {\portadaTitulo};
  \node[anchor=north, text=asgSlate, align=center, font=\sffamily\large]
    at ($(current page.north)+(0,-13.5)$) {\portadaSubtitulo};

  % Línea de acento
  \fill[asgAccent] ($(NO |- corte)+(0,0.12)$) rectangle ($(SE |- corte)+(0,-0.12)$);

  % Zona inferior clara
  \node[anchor=north, text=asgNavy, font=\sffamily\bfseries\fontsize{30}{34}\selectfont]
    at ($(corte)+(0,-2.3)$) {\portadaDocumento};
  \fill[asgAccent] ($(corte)+(-1.5,-4.0)$) rectangle ($(corte)+(1.5,-4.08)$);

  \node[anchor=north, align=center, font=\sffamily]
    at ($(corte)+(0,-5.2)$) {%
      \begin{tabular}{r@{\hspace{0.6cm}}l}
        {\color{asgSlate}\small AUTOR}    & {\color{asgNavy}\Large \portadaAutor} \\[0.45cm]
        {\color{asgSlate}\small DIRECTOR} & {\color{asgNavy}\Large \portadaDirector} \\
      \end{tabular}};

  \node[anchor=south, text=asgSlate, font=\sffamily]
    at ($(current page.south)+(0,1.6)$) {\portadaFecha};
\end{tikzpicture}
\end{titlepage}

% ÍNDICE
\frontmatter

% CONTROL DE VERSIONES
% Para una nueva versión, añadir una fila al final de la tabla.
\chapter*{Control de versiones}

\begin{tabularx}{\textwidth}{@{}l l l X@{}}
  \toprule
  \textbf{Versión} & \textbf{Fecha} & \textbf{Autor} & \textbf{Descripción} \\
  \midrule
  1.0 & 30/09/2026 & \portadaAutor & Versión inicial del manual de usuario. \\
  \bottomrule
\end{tabularx}

\tableofcontents
\listoffigures
\lstlistoflistings

% CONTENIDO
\mainmatter
\subfile{instalacion}
\subfile{manual}
\subfile{catalogo}

\end{document}
